\section{相关工作}

\paragraph{基准评测的科学。}
近期一系列工作将注意力从提高排行榜分数，转向审视基准分数本身是否提供可靠的科学证据。除数据污染之外——即基准测试题出现在模型训练数据中、从而可能抬高评测分数~\cite{oren2024proving,sun2025emperor,zhu2024dynamic}——已有研究表明，基准结果对解码、提示与实现细节敏感~\cite{hochlehnert2025sober}，并且对表面措辞或数据集缺陷十分脆弱~\cite{mousavi2026garbage}。另有工作从信号与噪声的角度研究评测可靠性~\cite{heineman2025signal}，或提出更广泛的基准质量准则，并表明广泛使用的基准在质量上差异很大~\cite{qian2026benchmark}。我们的工作补充了这些努力：将基准表述为难度上的经验测度，从而能够严格形式化测度失配与覆盖缺口。

\paragraph{IRT 及其应用。} IRT 起源于心理测量学，是教育与心理测量的一个框架~\cite{embretson2013item,lord2012applications}。在 NLP 中，\cite{lalor2016building} 将 IRT 引入系统评测，\cite{rodriguez2021evaluation} 将这一视角扩展到排行榜分析。此后，IRT 成为 LLM 基准中难度感知评测的广泛基础。Easy2Hard-Bench~\cite{ding2024easy2hard} 使用 IRT 为基准题目赋予标准化的连续难度标签；它还将这些标签与人类成对判断进行验证，并报告其一致性优于基于 GPT-4 的难度标注器。\cite{zhou2026lost} 使用基于 IRT 的框架分析现有 LLM 基准的测量质量。
% Beyond benchmark analysis, IRT has also been used for efficient evaluation~\cite{kipnis2024metabench,polo2024tinybenchmarks,zhang2025benchmark}, whose goal is to preserve the information of much larger benchmarks using far fewer evaluation items, and for adaptive testing, whose goal is to select the most informative items for each model based on its estimated ability during evaluation~\cite{hofmann2025fluid}.

% \paragraph{Rethinking benchmark analysis.} A recent line of work has shifted attention from improving leaderboard scores to examining whether benchmark scores themselves provide reliable scientific evidence. \cite{hochlehnert2025sober} shows that reasoning benchmark results can vary substantially under seemingly minor evaluation choices, such as decoding settings, prompt format, or implementation details, and calls for more standardized and reproducible evaluation. \cite{heineman2025signal} studies benchmark quality through the lens of signal, the ability to separate models, and noise, the instability induced by random variation, and analyzes practical ways to improve reliability. \cite{mousavi2026garbage} further shows that reasoning benchmark scores can be brittle to surface wording and dataset flaws, motivating more diagnostic and interpretable evaluation. \cite{qian2026benchmark} proposes explicit criteria for benchmark quality, including ranking consistency, discriminability, and capability alignment, and shows that widely used benchmarks vary substantially along these dimensions. Our work differs by formulating a benchmark as an empirical measure over the difficulty axis, which enables us to analyze how measure mismatch distorts aggregate score comparisons and how discrete coverage gaps hide localized model improvements.

